% ======================================================================
% SECTION 6: HR 4504 - Reduction of Human Doctor and Nurse Error Rate Act of 2026
% Adaption of HTI-1 DSI Final Rule 2023 and FDA AI Decision-Making
% Draft Guidance 2025.
% ======================================================================

[PROCESSING INSTRUCTION (HR 4504 - Reduction of Human Doctor and Nurse
Error Rate Act of 2026, 7 to 10 paragraphs structured into the five
legislative-act subsections, plus one ASCII comparison diagram):

The bill must explicitly state in its first paragraph that it is an
"Adaption of HTI-1 DSI Final Rule of 2023 and FDA AI Regulatory
Decision-Making Draft Guidance 2025" so that readers can trace the
chain of prior authority. Cite \cite{rank9_onc_hti1_2023},
\cite{hti1_dsi_fact_sheet_2023}, and
\cite{fda_ai_regulatory_decision_draft_2025}.

Subsection 6.1 - Findings and Declarations: This is the core bill for
the project brief's third aim ("reduce human doctor and nurse error
rates and lack of physical and intellectual capabilities"). Read the
following in full:
- new-trial/national-24-7-trial/paper/full-paper/sections/results.tex
  (the four LLM simulations demonstrating that AI patient prediction
  outperforms current narrow supervised models)
- new-trial/national-24-7-trial/Background-A/chunk\_01\_baseline\_and\_short\_horizon.md
  and chunk\_02\_multimodal\_and\_limitations.md (Manz 2020 AUC 0.89,
  Sidey-Gibbons 2022, the Huang 2025 null result of ML versus Cox)
- new-trial/national-24-7-trial/Background-B/chunk\_01\_baseline\_and\_prediction\_domains.md
  and chunk\_02\_response\_metrics\_conclusions.md (SCORPIO, PROGPATH,
  multimodal foundation models)
- patients/paper/Deep-Research-3/part2\_ai\_robotics\_legislation.md
  (the SHIELD-RT prospective RCT: 45 percent reduction in acute-care
  events, 48 percent reduction in costs)

Compose 5 to 7 enumerated findings citing
\cite{kawchak_2026_19994945} for the four-simulation paper,
\cite{vorontsov2024virchow}, \cite{rocque2025remote_symptom_monitoring},
\cite{gonzalez2025bayesian_urgent_care_nsclc},
\cite{mazor2025ai_notifications_trial}, and the FDA RTCT announcement
at \cite{fda_rtct_press_2026}. The findings must:
- Acknowledge respectfully that human doctors and nurses make errors
  at documented rates (cite the prior published rates rather than
  inflating them).
- Acknowledge that AI and robotics can reduce specific error
  categories where prospectively validated (SHIELD-RT, ASyMS,
  eRAPID).
- Acknowledge that humans retain irreplaceable roles in empathy,
  consent counseling, and ethical oversight, and that the bill does
  NOT eliminate human doctors or nurses; it shifts the patient's
  default for procedurally definable tasks to physical-AI execution
  where the patient prefers and the evidence supports.

Subsection 6.2 - Definitions:
- "Documented human error rate" means the per-task error rate
  published in peer-reviewed literature for the specific task
  category (surgical site infection, medication reconciliation,
  needle placement deviation, RT positioning, sepsis recognition,
  etc.).
- "Physical AI substitute" means a Physical AI system, robot, or
  agent that can perform the same task category with a documented
  error rate at or below the human baseline. Reference the 53 core
  agents at sponsor/final\_paper/scripts/core\_agents/ (in particular
  safety\_agent.py, clinical\_accountability\_agent.py, and
  audit\_trail\_manager.py) for the operational baseline.
- "Patient-initiated AI second opinion" means the patient's right to
  request, before any consequential clinical decision, a parallel
  evaluation by a qualified AI system meeting HTI-1 DSI transparency
  requirements.

Subsection 6.3 - Error Rate Reduction Rights (operative core):
- Right to a published per-site, per-task human error rate prior to
  consenting to any procedure.
- Right to substitute a Physical AI executor for any task where the
  documented Physical AI error rate is below the human baseline AND
  the patient prefers AI execution.
- Right to a patient-initiated AI second opinion, generated by a
  qualified system, prior to any consequential clinical decision
  (surgery, chemotherapy regimen change, end-of-life care decision).
- Right to mandatory error tracking across all sites, with public
  per-site reporting; reference the dashboard at
  sponsor/final\_paper/scripts/dashboard/ and the 168-hour cumulative
  metrics at sponsor/final\_paper/168\_hours/instructions/core\_i5\_6200u\_4gb/sponsor\_168h\_summary.json
  for the operational reporting pattern.
- Right to refuse a human-only execution path for any task where a
  qualified Physical AI substitute exists, mirroring the FDA RTCT
  pilot's "real-time signal sharing" model. Cite
  \cite{fda_rtct_press_2026}.

After Subsection 6.3, embed a comprehensive ASCII comparison diagram
in a verbatim block. Three columns (Task Category, Human Baseline
Error Rate, Physical AI Error Rate) and rows for: surgical site
infection, medication reconciliation, needle placement deviation, RT
positioning, sepsis recognition, AE detection from ePRO. Source the
human baselines from the prior literature cited in
patients/paper/Deep-Research-1/chunk\_1\_foundations.md and the AI
baselines from
new-trial/national-24-7-trial/Background-A/ and
sponsor/final\_paper/scripts/output/sponsor\_24h\_summary.json.

Subsection 6.4 - Implementation: 18-month deadline for ALL
CMS-billable sites to publish per-task human and AI error rates;
HTI-1 DSI transparency requirements apply to every AI second-opinion
system; FDA, AHRQ, ONC, and CMS roles. Reference the patient priority
operational baseline at
new-trial/national-24-7-trial/paper/full-paper/sections/discussion.tex.

Subsection 6.5 - Reporting and Enforcement: Annual federal error-rate
report by site, by task, by AI substitute (if any); CMS
reimbursement penalty for sites that fail to publish error rates;
information-blocking penalty harmonized with
\cite{federal_register_info_blocking_disincentives_2024}.

Close the bill (1 paragraph after Subsection 6.5) with two patient-
control framings:
(i) FOR INDIVIDUAL PATIENTS: A 58-year-old female NSCLC patient
    matching PAT-2026-0042 reviews the published per-site error
    rates at her chosen site, requests an AI second opinion before
    her cycle-1 pembrolizumab dosing decision, and refuses a
    human-only RT positioning step in favor of the RTPOS-01 robot
    documented at hour-12 of
    new-trial/national-24-7-trial/hour-12/robot\_logs.md.
(ii) FOR BROAD ADOPTION: All U.S. CMS-billable oncology trial sites
     must publish per-task error rates within 18 months. The bill
     positions the United States as Number 1 in the world for
     evidence-based patient-led error reduction. The bill explicitly
     respects the human doctor and nurse roles that remain
     irreplaceable (empathy, consent counseling, ethical oversight)
     and constrains the AI-substitution path to procedurally
     definable tasks with prospectively validated error reductions.
     Cite \cite{kawchak_2026_19994945} and \cite{kawchak_2026_19119939}.]
